\scanpage{13}
\handwritten{\mspage{4}}
b)\quad Pour toute courbe algébrique (propre, lisse, connexe) $\unsure{\mathbb{C}|C}$ sur $k$, et tout cycle algébrique $Z$ sur $Y\mathbin{\mathrm{x}}C$ de codimension $n$, l'homomorphisme correspondant
\[
Z:H^1(C)\underline{\deleted{Z}\handwritten{Q}}_{\unsure{\chi|\ell}}(1))\handwritten{\longrightarrow}H^{2n-1}(Y,\underline{Q}_{\chi}(n))
\]
a son image dans celle de $i^{\handwritten{*}}$ de (***). (NB Si on admet le yoga des motifs, justifiable par les``conjectures standard'', cetté\deleted{\unsure{xxxxé|même}} condition signifie que le \inserted{\handwritten{plus grand}} morceau \deleted{\unsure{dx|du|de}} ``de niveau 1'' du motif \deleted{Coker} ayant comme réalisation $\chi$-adique $\mathrm{Coker}\ i^{\handwritten{*}}$ est nul).

\uline{Sous ces conditions}, \uline{pour tout cycle algébrique} $z\handwritten{\in}Z^n(Y)$ \uline{qui est algébriquement équivalent à zéro}, \uline{il existe un cycle} $z'\handwritten{\in}Z^n(X)$ \uline{algébriquement équivalent à zéro}, \unsure{\uline{et un}|et \uline{un}} $z''\handwritten{\in}Z^n(Y)$, tels que
\[
z=i^{\handwritten{*}}(z')+z''\tag*{a')}
\]
b')\quad \deleted{\unsure{l'homx|l'hom}} $z''=Z^{\handwritten{''}}(t_1)-Z^{\handwritten{''}}(t_{\unsure{o|0}})$, où $Z''$ est un \inserted{\unsure{$\underline{Q}$-|\handwritten{$\underline{Q}$-}}}cycle \uline{de codimension} $n$ \uline{sur un}\deleted{e courbe} \uline{produit} $\deleted{C\mathbin{\mathrm{x}}}\deleted{X}\handwritten{Y}\mathbin{\mathrm{x}}C$, \uline{induisant un homomorphisme} \uuline{nul} \deleted{de} $H^1(C,\underline{\deleted{Z}\handwritten{Q}}_{\chi}(1))\handwritten{\longrightarrow}H^{2n-1}(Y,\underline{Q}_{\chi}(n))$.

Démonstration immédiate : on écrit $z=Z(t_1)-Z(t_o)$ avec $Z\handwritten{\in}Z^n(Y\mathbin{\mathrm{x}}C)$, ce $Z$ définit un homomorphisme $Z:H^1(C)\handwritten{\longrightarrow}H^{2n-1}(Y)$, \deleted{dont le} qui par hypothèse b) tombe dans $\mathrm{Im}\ i^{\handwritten{*}}$ donc par hypthèse a) est égal au composé $(iv)Z=i(vZ)$, or $vZ:H^1(C)\handwritten{\longrightarrow}H^{2n-1}(X)$ est algébrique puisque $v$ l'est, donc donné par un cycle algébrique $Z'$ sur $X\mathbin{\mathrm{x}}C$, et si on pose $Z''=Z-(i\mathbin{\mathrm{x}}\mathrm{id}_C)^{\handwritten{*}}(Z\handwritten{'})$, on aura par construction que $Z''$ induit l'homomorphisme nul $H^1(C)\handwritten{\longrightarrow}H^{2n-1}(Y)$. \deleted{D'autre part} Il suffit alors de poser $\unsure{z'|Z'|\mathbb{Z}'}=Z'(t_1)-Z'(t_o)$, $z''=Z''(t_1)-Z''(t_o)$.

3.\quad Mettant ensemble les réflexions de 1) et 2), on trouve ceci. On se donne des schémas propres et lisses $X$ et $Y$ sur $S$ \inserted{\handwritten{\unsure{règulier|régulier}}} connexe, et un $S$-morphisme
\[
i:Y\handwritten{\longrightarrow}X\ .
\]
% archive sheet number: 12
